% ECONOMICS_PROBLEM: {"id": "D82-189", "title": "Voluntary participation over time", "jel_code": "D82", "source_id": "OP-0291"}
% !TeX program = xelatex
\documentclass[11pt,a4paper]{article}
\usepackage[margin=23mm]{geometry}
\usepackage{fontspec}
\setmainfont{Latin Modern Roman}
\usepackage{amsmath,amssymb,mathtools}
\usepackage{xcolor,enumitem,booktabs,longtable,array,xurl,hyperref}
\definecolor{muted}{HTML}{53636C}
\hypersetup{colorlinks=true,urlcolor=blue,linkcolor=blue}
\setlength{\parindent}{0pt}
\setlength{\parskip}{5pt}
\newcommand{\R}{\mathbb R}
\newcommand{\E}{\mathbb E}
\newcommand{\N}{\mathbb N}
\newcommand{\ind}{\mathbf 1}
\newcommand{\Var}{\operatorname{Var}}
\newcommand{\Cov}{\operatorname{Cov}}
\newcommand{\supp}{\operatorname{supp}}
\DeclareMathOperator*{\argmin}{arg\,min}
\DeclareMathOperator*{\argmax}{arg\,max}
\newcommand{\entrymeta}[3]{{\sffamily\small\color{muted}\textbf{\texttt{#1}}\quad|\quad #2\par #3\par}}
\newcommand{\Problemref}[1]{\texttt{#1}}
\newcommand{\JELref}[2]{\texttt{#2} (JEL \texttt{#1})}
\begin{document}
\section*{Voluntary participation over time}
\textbf{Problem ID:} \texttt{D82-189}\quad \textbf{JEL:} \texttt{D82}\par
% BEGIN ECONOMICS STATEMENT
\label{problem:OP-0291}
{\sffamily\small\color{muted}Original subject group: Mechanism design and market design\par}
\label{definitions:problem:30007573}
\textbf{Audit classification:} Published research agenda. The survey’s §8 discusses participation and initial bonding. Wealth-limited escrow design is an editorial extension; finite-history optimization alone is not a new existence problem.
\subsection*{Definitions and precise research target}
\textbf{Model and research target.} Let one seller repeatedly allocate a service to \(n\) agents over \(T\) periods. Agent \(i\)'s type follows a known finite Markov chain and yields quasilinear discounted utility, with common discount factor \(\beta\in(0,1]\). A direct mechanism chooses feasible allocation \(x_t\) and nonnegative payment \(p_{it}\) from report histories. Agent \(i\) has observable initial liquid wealth \(B_i\), cannot borrow, and may exit after any private history with no future allocations or payments. Impose the pathwise constraint \(\sum_{s\leq t}p_{is}\leq B_i\), dynamic incentive compatibility against every reporting-and-exit strategy, and sequential participation
\[
 E\left[\sum_{s=t}^{T}\beta^{s-t}
 (v_i(x_s,\theta_{is})-p_{is})\mid h_{it}\right]\geq0
\]
at every truthful private history \(h_{it}\). Characterize revenue-optimal mechanisms and the loss relative to mechanisms requiring only initial participation, as functions of wealth, type persistence, and horizon. Determine whether limited refundable escrow or capped deposits recover part of that loss, stating the cash and exit rules exactly before comparing mechanisms. An existence or characterization result must include off-path incentive restrictions and histories reached after other agents' reports.

\emph{Definitions and maintained assumptions (editorial unless a source definition is named).} The known joint initial and Markov type law, allocation feasibility and valuation functions are primitives; agents know their own types and the specified public information. Exit at the beginning of period \(t\) eliminates future trades and payments and gives normalized future utility zero; previously received service and paid fees remain sunk. With nonnegative payments and no wealth replenishment, \(\sum_{s\le t}p_{is}\le B_i\) is a pathwise liquidity constraint in nominal currency, distinct from discounted utility. Dynamic IC compares truth-and-continuation with every nonanticipating reporting and exit strategy at every feasible private history, using a specified off-path posterior rule or ex-post robustness. For finite types/horizon, randomized mechanisms with bounded payments can be encoded by finitely many history variables and linear IC/IR/liquidity constraints; mere numerical computability is not the open issue. Seek structural formats, bounds uniform in horizon and polynomial representation/computation under explicitly encoded primitives. Refundable escrow needs a cash-balance process recording deposits, interest, forfeiture and refunds; it cannot simply be added to a no-refund cumulative-payment constraint.
\subsection*{Variants and assumption changes}
\begin{enumerate}
\item \textbf{No escrow (editorial).} Study a single buyer with two Markov types, nonnegative fees, zero future exit payoff and initial wealth \(B\). Characterize horizon-uniform structural policies and the revenue gap from initial-only participation.
\item \textbf{Refundable deposit (editorial).} Permit an initial deposit at most \(B\), immediate refunds on exit and specified deductions. Compare revenue under the same cash-flow feasibility and truthful-exit constraints; unrestricted bonding is a different model.
\end{enumerate}
\subsection*{References and source audit}
\begin{enumerate}
\item §8 printed p.50 in June2018 revision. Supports the stated research area; exact displayed specification is editorial. DOI publication is2019; inspected revision is June2018. \url{https://cowles.yale.edu/sites/default/files/2022-09/d2102-r.pdf}
\end{enumerate}
\begin{enumerate}\raggedright
\item\hypertarget{definitions:source:30007573:1}{} Dirk Bergemann; Juuso Välimäki. 2019. \emph{Dynamic Mechanism Design: An Introduction}. Published JEL57(2),235–274 (2019); inspected June19,2018 Cowles revision §8, printed pp.49–51.
\url{https://cowles.yale.edu/sites/default/files/2022-09/d2102-r.pdf}.
DOI: \url{https://doi.org/10.1257/jel.20180892}.
\end{enumerate}

\subsection*{Common conventions: Source mathematical conventions}

These conventions apply to each entry unless explicitly replaced.

\textbf{Models and identification.} A primitive model \(m\) specifies all probability laws, feasible choices, information, equilibrium and measurement rules used by its target. The admissible class \(\mathcal M\) is fixed independently of outcomes. If \(P_O(m)\) is its observable probability law and \(\tau(m)\) the target, its identified set at observable law \(P\) is
\[
 \mathcal I_\tau(P)=\{\tau(m):m\in\mathcal M,\ P_O(m)=P\}.
\]
Point identification means that this set is a singleton. An empty set signals incompatibility of data and assumptions. Coordinate bounds need not describe the joint set. Bounds are sharp only if every retained target is attainable or arbitrarily approachable by compatible models. Finite bounds require explicit restrictions on unobserved quantities; unrestricted confounding, hidden wealth, extrapolation or arbitrary equilibrium selection can yield uninformative sets.

\textbf{Causal interventions.} A potential outcome \(Y_i(p)\) is the outcome under a complete policy regime \(p\), including timing, financing, information and market spillovers. The target population and units are fixed before treatment. With interference, use the complete assignment vector or a justified exposure mapping; individual no-interference notation alone cannot identify an economy-wide intervention. A difference of mean potential outcomes does not require the cross-world joint distribution. Conditioning on post-treatment migration, survival, enrollment or employment changes the estimand and needs separate justification.

\textbf{Statistical statements.} An estimator, confidence set, forecast or test specifies a sampling law, model class and loss/coverage criterion. Uniform \((1-\alpha)\)-coverage means \(\inf_{m\in\mathcal M}\Pr_m\{\tau(m)\in C_n\}\geq1-\alpha\), or an explicitly asymptotic version. The whole parameter space trivially has coverage, so informativeness requires a stated width, risk or power objective. A forecasting improvement is relative to a named benchmark, information set and evaluation population.

\textbf{Equilibrium and optimization.} An equilibrium comprises feasible plans, optimality, beliefs where needed, budget constraints and all market/resource clearing. The primitives or a stated benchmark define each correspondence \(\mathcal E(p;m)\). Nonemptiness is assumed or proved, never inferred just from compact policy sets. A selected-equilibrium target uses a declared selection rule; a robust target quantifies over all permitted equilibria. Use suprema or infima unless attainment is established. An \(\varepsilon\)-optimal policy is feasible and within \(\varepsilon\) of the finite optimum value. Report an empty feasible set separately.

\textbf{Welfare and accounting.} Normative utility, interpersonal normalization, social weights, numeraire, discounting and the use of public receipts are primitives. Behavioral utility need not coincide with the welfare criterion. Household welfare includes ownership income and rents once; adding profits again to compensating variation generally double-counts them. A monetary equivalent is defined by an explicit transfer and price convention, with monotonicity and range conditions. Average effects, mechanism contributions, transfers and welfare losses are different targets. Infinite sums require convergence and dynamic feasible plans require terminal or no-Ponzi conditions.

\textbf{Source versions.} A working-paper date, revision date and journal publication year are distinguished. A DOI for a journal version is not automatically the DOI of an earlier manuscript. Locators refer to the stated version and printed pages where specified. A metadata-only check does not verify a claim about a full-text conclusion. Every entry's source subsection narrows the provenance claim accordingly.
% END ECONOMICS STATEMENT
\end{document}
